\section{Modelos UML y de procesos}
\label{sec:uml}

\subsection{Cobertura del modelado}

\begin{table}[H]
\centering\footnotesize
\caption{Modelos elaborados y su función en la especificación}
\label{tab:modelos}
\begin{tabular}{|L{4.2cm}|C{1.6cm}|L{8.8cm}|}
\hline
\rowcolor{grisclaro}\textbf{Modelo} & \textbf{Cantidad} & \textbf{Qué responde que los demás no responden}\\ \hline
Diagrama de contexto & 1 & Dónde está la frontera del sistema y qué queda fuera\\ \hline
Modelado i* (SD y SR) & 2 & Qué depende cada actor de los demás, y por qué querría usar el sistema\\ \hline
Casos de uso & 18 & Qué secuencia de pasos produce cada resultado, con sus alternativas y excepciones\\ \hline
Diagrama de clases refinado & 1 & Qué entidades del dominio existen y cómo se relacionan\\ \hline
Diagramas de secuencia & 3 & Cómo colaboran los objetos en los flujos más críticos\\ \hline
Diagrama de actividad & 1 & Cómo se encadena el ciclo semanal completo\\ \hline
Diagramas de estados & 2 & Qué ciclo de vida siguen el racimo y la alerta\\ \hline
Diagramas de flujo de datos & 2 & Qué proceso transforma qué dato y dónde se almacena\\ \hline
Componentes y despliegue & 2 & Cómo se reparte la funcionalidad y dónde se ejecuta\\ \hline
\end{tabular}
\end{table}

\subsection{Lectura interpretativa de los modelos principales}

\subsubsection*{Casos de uso: el hallazgo de la auditoría}

El diagrama general identifica dieciocho casos de uso, pero hasta esta unidad solo diez tenían
especificación textual. Los ocho restantes ---\id{CU-11} a \id{CU-18}--- existían como nodo del
diagrama y como referencia en la matriz, sin flujo, sin precondición y sin excepciones.

El vacío no lo detectó ninguna revisión de contenido: lo detectó la métrica M1b al dividir casos
especificados entre casos identificados y obtener 55,6\,\%. Es un ejemplo de por qué una auditoría
cuantitativa aporta algo que la lectura no aporta: un lector que recorre el documento de principio a
fin ve diez casos de uso bien especificados y no echa en falta los que no están.

Los ocho se especificaron en esta unidad con el mismo formato que los diez anteriores. Dos merecen
comentario. \id{CU-16} (generar alerta temprana) tiene como actor primario al propio sistema, que
actúa sin intervención humana; su flujo alternativo B evita la duplicación de alertas equivalentes,
que es el defecto que produciría un panel inservible por saturación. \id{CU-18} (explicar una
decisión automática) es incluido obligatoriamente por los tres casos de uso con decisión automática,
de modo que la explicabilidad no es una función que el usuario invoque sino una condición de la
propia emisión del resultado.

\subsubsection*{Diagramas de estados: dos entidades y por qué solo dos}

Se modelan el racimo y la alerta. La elección no es arbitraria ni exhaustiva: son las dos entidades
del dominio cuyo ciclo de vida determina una consecuencia económica y cuya transición no es trivial.

El estado \emph{Degradado} del racimo materializa el hallazgo de \id{EV-04} sobre la ventana entre
corte y entrega: el material híbrido tolera un intervalo mayor por su contenido de antioxidantes,
mientras que el \emph{guineensis} se deshidrata en dos o tres días. Sin ese estado, la ventana
corte--entrega de \id{RF-25} sería una regla sin objeto sobre el que aplicarse.

El estado \emph{Descartada} de la alerta corresponde a los falsos positivos del modelo de diagnóstico
y alimenta el ciclo de reentrenamiento previsto en \id{RNF-15} y \id{RNF-17}. Es el punto donde el
modelo de estados y los componentes de inteligencia artificial se tocan: un falso positivo no es solo
un error que se descarta, es un dato de entrenamiento.

\subsubsection*{Diagramas de flujo de datos: por qué se incorporaron ahora}

Los diagramas de flujo de datos se incorporaron en esta unidad, no en la de modelado. La razón es
concreta: la cadena de trazabilidad extremo a extremo exige enlazar cada requisito con el proceso que
lo realiza, y sin diagrama de flujo de datos esa columna quedaría vacía.

La decisión de diseño que conviene declarar es la correspondencia uno a uno entre los siete procesos
del nivel 1 y los siete bloques funcionales de la \S3.2 del ERS. No es una coincidencia del dominio:
se adoptó deliberadamente para que la introducción del modelo no obligara a reclasificar ningún
requisito ya especificado y para que un cambio en un requisito se propague por una sola ruta. La
alternativa ---descomponer el sistema por criterio propio del análisis de flujo--- habría producido
un modelo posiblemente más fiel pero con dos taxonomías paralelas del mismo sistema, que es
exactamente el mecanismo que generó los defectos de consistencia detectados en la Unidad IV.

\subsubsection*{Modelado i*: lo que aportó y lo que no}

Los diagramas SD y SR modelan las dependencias entre actores antes de que exista un sistema. Su
aporte concreto al proyecto fue detectar que la planta extractora es una parte interesada con poder
alto e interés medio, cuya calificación de calidad es una dependencia externa que el sistema no
controla pero de la que depende una función completa (\id{RF-23}, \id{RF-24}).

Lo que no aportó, y conviene decirlo: el modelado i* no produjo requisitos que la elicitación no
hubiera producido ya. Su valor fue de encuadre y de verificación cruzada, no de descubrimiento.

